\documentclass[lualatex]{ltjsarticle}
\usepackage{luatexja}
\usepackage{amsmath, amssymb, amsthm}
\usepackage{mathtools}
\usepackage{tikz-cd}
\usepackage{physics}
\usepackage{hyperref}
\hypersetup{
  colorlinks=true,
  linkcolor=blue,
  citecolor=blue,
  urlcolor=blue
}

\title{行列の階数と退化次数に関するメモ}
\author{CODEX (OpenAI)\thanks{本ノートは CODEX (OpenAI) により作成された TeX ドキュメントです。} \\ CODEX (OpenAI)\thanks{付随する Lean ファイルも CODEX (OpenAI) により編集されました。}}
\date{2025年9月19日}

\newtheorem{theorem}{定理}
\newtheorem{lemma}[theorem]{補題}
\newtheorem{proposition}[theorem]{命題}
\theoremstyle{definition}
\newtheorem{definition}[theorem]{定義}
\theoremstyle{remark}
\newtheorem{remark}[theorem]{備考}

\begin{document}

\maketitle

\section{概要}
本ノートでは、学習課題で扱った具体的な実行列 $A, B, C, D, E$ を題材に、階数 (rank) と退化次数 (nullity) の基本性質を整理する。特に、Lean での形式化で得られた知見を踏まえ、手計算と抽象的議論の橋渡しを意図した解説を行う。

\section{定義の整理}
\begin{definition}[列の線形独立性]
$m\times n$ 行列 $M$ の列ベクトル集合から有限部分集合 $S \subseteq \{0,\dots,n-1\}$ を選ぶ。係数列 $c = (c_j)_{j \in S}$ に対し、以下が成り立つとき $S$ は線形独立という。
\[
  \sum_{j\in S} c_j M_{\bullet j} = 0 \implies \forall j\in S,\; c_j = 0.
\]
\end{definition}

\begin{definition}[階数と退化次数]
列ベクトルの線形独立な最大個数を階数 $\operatorname{rank}(M)$ とし、列数を $n$ とするとき退化次数を $\operatorname{nullity}(M) = n - \operatorname{rank}(M)$ と定める。
\end{definition}

\section{具体例の検討}
\subsection{上三角行列 $B$}
行列
\[
B = \begin{pmatrix}
1 & 2 & 3 \\
0 & 1 & 2 \\
0 & 0 & 1
\end{pmatrix}
\]
は上三角かつ対角要素が全て $1$ であるため、列（あるいは行）が完全に独立であり、\(\operatorname{rank}(B) = 3\) である。Lean では `norm_num` を用いて $B_{ii} \neq 0$ を確認することで形式的な証明が成立した。

\subsection{行列 $C$ と rank-nullity}
行列
\[
C = \begin{pmatrix}
1 & 0 & 2 \\
0 & 1 & -1
\end{pmatrix}
\]
に対して、第1列・第2列が独立であることを成分比較によって確認できる。従って $\operatorname{rank}(C)=2$。列数は 3 なので $\operatorname{nullity}(C)=1$。これは後述する rank-nullity 定理の具体例である。

\subsection{行列 $D$ の行依存性}
\[
D = \begin{pmatrix}
1 & 2 & 3 \\
2 & 4 & 6 \\
0 & 1 & 1 \\
1 & 3 & 4
\end{pmatrix}
\]
第2行は第1行の2倍であるため、行集合に線形従属が存在する。Lean では `ext`/`fin_cases` による成分比較と `norm_num` を組み合わせて `D.row 1 = 2 • D.row 0` を確定させた。

\section{階数と行列式}
行列 $E$ は
\[
E = \begin{pmatrix}
1 & 2 & 3 \\
4 & 5 & 6 \\
7 & 8 & 9
\end{pmatrix}
\]
であり、直接展開により行列式が $0$ である。これは $E$ の階数が 3 未満であることを意味する。Lean 上では定義式を `simp` と `norm_num` で逐一評価し、零であることを形式的に示した。

\section{図式による視覚化}
列空間と核の関係を図式で表すため、線形写像 $T : \mathbb{R}^3 \to \mathbb{R}^2$ を $C$ と対応付けて考える。
\[
\begin{tikzcd}
0 \arrow[r] & \ker T \arrow[r, hook, "\iota"] & \mathbb{R}^3 \arrow[r, "T"] & \operatorname{im} T \arrow[r] & 0
\end{tikzcd}
\]
この正確列は $\dim \ker T + \dim \operatorname{im} T = 3$ を与え、rank-nullity 定理そのものを表現する。$C$ の場合、$\dim \operatorname{im} T = 2$ なので、核の次元は $1$ となる。

\section{Lean 形式化のポイント}
\begin{itemize}
  \item `Finset.sum` を用いることで、列（行）の線形結合を有限集合上の総和として扱った。
  \item `classical` ブロックを導入し、`Finset.max'` に必要な選択公理を活用して最大値を取得した。
  \item `norm_num` や `simp` を駆使して具体的数値計算を自動化し、手計算の誤りを排除した。
\end{itemize}

\section{まとめ}
Lean による形式検証を通じて、
\begin{itemize}
  \item 上三角行列の対角要素非零性 \Rightarrow フルランク、
  \item 列の独立性による階数計算、
  \item 行の比例関係による従属性の検出、
  \item 行列式と階数の関係、
  \item rank-nullity 定理の具体例
\end{itemize}
を明確に確認した。これらの結果は、行基本変形や核・像の概念を可視化し、線形代数の基礎的な理解を支援するものである。

\end{document}
